\begin{ExerciseBlock}
\ExerciseHeading{06}{Abbremsen eines Zylinders}{rigid_bodies_braking_a_hollow_cylinder}
In der Abbildung ist ein rotierender Eisenzylinder (Dichte $\rho_{Fe} = 7,874\,$g/cm³) mit einem Radius $R=20\,$cm und einer Höeh $h=10\,$cm zu sehen. Dieser rotiert mit einer Frequenz $f_0=100\,s^{-1}$. Anschließend wird eine Bremsscheibe mit der Kraft $F_N$ auf den Zylinder gepresst. Der Gleitreibungskoeffizient ist $\mu=0,2$.

\begin{enumerate}
	\item Bestimmen Sie die Winkelgeschwindigkeit $\omega$ der Rotation.
	\item Wie groß muss die Kraft $F_N$ sein, damit der Zylinder innerhalb einer Sekunde zum stehen kommt?
	\item Wie viele Umdrehungen vollzieht der Zylinder bis dieser stillsteht?
\end{enumerate}

\begin{center}
	\CachedTikzFigure{content/resources/chapters/rigid_bodies/tikz/braking_a_hollow_cylinder.tex}
\end{center}
\end{ExerciseBlock}
\UseExerciseSolution{rigid_bodies_braking_a_hollow_cylinder}
